\documentclass[journal]{IEEEtran}

\usepackage[utf8]{inputenc}
\usepackage[T1]{fontenc}
\usepackage{amsmath,amssymb}
\usepackage{graphicx}
\graphicspath{{../Versions-html-and-docx/media/}}
\usepackage{booktabs}
\usepackage{multirow}
\usepackage{array}
\usepackage{url}
\usepackage{cite}
\usepackage{siunitx}
\usepackage{xcolor}
\usepackage{balance}

\sisetup{
  group-separator = {\,},
  detect-all
}

\newcommand{\Fp}{F_\mathrm{p}}
\newcommand{\Fn}[1]{F_{#1}}
\newcommand{\Hz}{\,\text{Hz}}
\newcommand{\dB}{\,\text{dB}}

\begin{document}

\title{Formant Analysis of Orchestral Instruments:\\A Spectral Reference for Timbral Convergence\\and Classical Doubling Practices}

\author{Yan~Maresz%
\thanks{Y.~Maresz is a composer and researcher. Correspondence: \texttt{y.maresz@cnsmdp.fr}}%
}

\markboth{IEEE/ACM Transactions on Audio, Speech, and Language Processing}%
{Maresz: Formant Analysis of Orchestral Instruments}

\maketitle

% ============================================================
\begin{abstract}
This paper presents the most comprehensive quantitative study to date of spectral formants across 30~orchestral instruments, analyzing \num{5914}~professionally recorded sustained samples from the SOL2020 (IRCAM) and Yan\_Adds databases. We introduce the \emph{Principal Formant} ($\Fp$), an energy-weighted spectral centroid that provides dramatically more stable timbral characterization than conventional peak-based formant extraction---reducing variance by up to $10.4\times$ for wide-tessitura instruments such as the trumpet ($\sigma_{\Fp}=\SI{98}{\hertz}$ vs.\ $\sigma_{\Fn{2}}=\SI{1018}{\hertz}$). Our central finding is a \emph{cluster of convergence} in the vowel zone /o/--/\aa/ (\SIrange{377}{506}{\hertz}), where 11~instruments from four orchestral families converge within \SI{129}{\hertz}. This convergence provides a measurable acoustic explanation for historically admired doublings: bassoon--violin ($\Delta=\SI{11}{\hertz}$), horn--viola ($\Delta=\SI{11}{\hertz}$), English horn--B$\flat$ clarinet ($\Delta=\SI{11}{\hertz}$). Cross-validation against four independent academic sources (Meyer, Backus, Giesler, McCarty/CCRMA) yields a 93\% agreement rate across 29~comparable instruments. Six quantitative orchestration principles are derived from the data, covering formant convergence, spectral complementarity, the ensemble section effect, mute-induced timbral transposition, cross-family formant groupings, and spectral masking avoidance.
\end{abstract}

\begin{IEEEkeywords}
Orchestral timbre, spectral formants, energy-weighted centroid, timbral convergence, orchestration, instrument doubling, SOL2020, spectral analysis.
\end{IEEEkeywords}

% ============================================================
\section{Introduction}
\label{sec:intro}

\IEEEPARstart{T}{he} timbral quality of orchestral instruments has been studied extensively through spectral analysis since the pioneering work of Helmholtz. While the formant structure of the human voice is well established through the source-filter model~\cite{fant1960}, the application of formant-based reasoning to orchestral instruments remains largely qualitative in orchestration treatises. Meyer~\cite{meyer2009} established vowel-like timbral correspondences for several instruments; Giesler~\cite{giesler1985} provided explicit vowel associations for 13~instruments; and Backus~\cite{backus1969} documented formant regions for brass and woodwinds. However, no prior study has systematically extracted formant data across the full orchestral palette from a unified corpus using a validated pipeline, nor has the acoustic basis of classical doublings been quantified through formant proximity.

The present work addresses this gap by analyzing \num{5914}~sustained samples covering 30~instruments, combining two spectral databases (SOL2020/IRCAM and the supplementary Yan\_Adds corpus). We introduce $\Fp$, an energy-weighted spectral centroid that overcomes the well-known instability of peak-based formant extraction for wide-tessitura instruments. The resulting dataset enables, for the first time, a quantitative framework for understanding orchestral doublings through measurable spectral convergence.

Our central discovery is a \emph{cluster of convergence} in the /o/--/\aa/ vowel zone (\SIrange{377}{506}{\hertz}), where instruments from four different families---strings, brass, woodwinds, and saxophones---share nearly identical first formants. This cluster provides the acoustic explanation for doublings that orchestrators have considered effective for centuries: bassoon with violin, horn with viola, English horn with B$\flat$~clarinet.

\subsection{Terminological Note}

Throughout this paper, $\Fn{1}$--$\Fn{6}$ denote \emph{spectral peaks} extracted by prominence-based peak detection on the spectral envelope, not vocal-tract formants in the strict sense of Fant~\cite{fant1960}. For cylindrical-bore instruments (clarinet, flute) where no fixed formant exists, $\Fn{1}$ designates the \emph{dominant spectral energy zone}. The notation is retained for consistency with the orchestration literature (Meyer, Giesler), where instrumental ``formants'' are standard terminology. The $\Fp$ measure (Section~\ref{sec:fp}) is explicitly defined as a spectral centroid and is not subject to this ambiguity.

% ============================================================
\section{Corpus and Data Sources}
\label{sec:corpus}

\subsection{Primary Databases}

The analysis draws on two complementary sample databases:

\begin{itemize}
\item \textbf{SOL2020 (IRCAM/EIC):} 16~solo instruments, \num{2391}~sustained samples with pre-computed spectral envelopes (specenv, \num{1024}~bins). Instruments include all standard woodwinds, brass, and strings in multiple extended techniques.
\item \textbf{Yan\_Adds:} 18~additional instruments, \num{1646}~samples sourced from VSL, Con~Timbre, and personal recordings. Covers string ensembles, auxiliary woodwinds (bass flute, contrabass flute, E$\flat$~clarinet, contrabass clarinet), and low brass (bass trombone, contrabass tuba).
\end{itemize}

\subsection{Academic Reference Sources}

Four independent sources serve as validation benchmarks:
\begin{enumerate}
\item Backus (1969)~\cite{backus1969}: historical reference, targeted measurements on brass and woodwinds.
\item Giesler (1985)~\cite{giesler1985}: detailed measurements with explicit vowel associations for 13~instruments.
\item Meyer (2009, 5th ed.)~\cite{meyer2009}: acoustic analysis with vowel-color zones and the \emph{Hauptformant} concept for strings.
\item McCarty/CCRMA (2003)~\cite{mccarty2003}: vowel space chart for orchestral instruments---used as directional reference only, as the author has acknowledged accuracy limitations.
\end{enumerate}

\subsection{Measurement Scope}

All formant values reported in this paper are measured on \textbf{sustained tones (ordinario)} only, excluding attack transients, extended techniques, and extreme dynamics. Muted variants are analyzed separately. Included techniques: normal, arco, sustain, vibrato, non-vibrato, flautando---dynamics \textit{pp} to \textit{ff}. Excluded: pizzicato, col legno, tremolo, trills, artificial harmonics, multiphonics, flutter tongue, glissandi.

% ============================================================
\section{Methodology}
\label{sec:method}

\subsection{Spectral Envelope Extraction}

\begin{table}[!t]
\centering
\caption{Pipeline Parameters}
\label{tab:pipeline}
\begin{tabular}{@{}ll@{}}
\toprule
Parameter & Value \\
\midrule
FFT size & \num{4096} samples \\
Sample rate & \SI{44100}{\hertz} \\
Frequency resolution & $\sim$\SI{10.8}{\hertz}/bin \\
Analysis range & \SIrange{100}{3000}{\hertz} \\
Input data & Pre-computed spectral envelopes \\
 & (specenv, \num{1024} bins) \\
Windowing/smoothing & Hann + smoothing (IRCAM \\
 & SuperVP/AudioSculpt upstream) \\
Peak detection & Most prominent peak \\
 & (not first peak) \\
Amplitude threshold & $-30\dB$ below regional maximum \\
Minimum distance & \SI{70}{\hertz} ($\sim$6.5 bins) \\
\bottomrule
\end{tabular}
\end{table}

The extraction pipeline (v22, validated) operates on pre-computed spectral envelopes from the IRCAM analysis chain (SuperVP/AudioSculpt), which applies Hann windowing and smoothing upstream. Peak detection uses \texttt{scipy.signal.find\_peaks} with prominence-based selection: the most \emph{prominent} peak is selected as $\Fn{1}$, not the first peak in frequency order. This methodological choice captures the perceptual principal formant, avoids secondary low-frequency resonances, and is consistent with auditory perception. Table~\ref{tab:pipeline} summarizes the pipeline parameters.

Two critical bugs were identified and corrected in earlier pipeline versions: (1)~rejection of samples with $\max(\mathrm{amplitude}) \leq 0$, which had eliminated 66.5\% of violin samples; (2)~linear division on dB values instead of dB subtraction for threshold computation. The corrected pipeline (v22) reproduces all 16~SOL2020 CSV instruments at $\Delta = \SI{0}{\hertz}$ versus raw specenv data.

\subsection{The Principal Formant ($\Fp$): Energy-Weighted Spectral Centroid}
\label{sec:fp}

For wide-tessitura instruments, the first spectral peak $\Fn{1}$ is highly register-dependent and does not characterize the instrument's global timbre. The violin, for example, exhibits $\sigma_{\Fn{1}} = \SI{235}{\hertz}$ across its range, while its $\Fn{2}$ shows $\sigma = \SI{651}{\hertz}$---a variation exceeding one octave.

We define the \emph{Principal Formant} $\Fp$ as the energy-weighted spectral centroid within an instrument-specific frequency band:
\begin{equation}
\Fp = \frac{\sum_{i} f_i \cdot A_i}{\sum_{i} A_i}
\label{eq:fp}
\end{equation}
where $f_i$ is the frequency of bin~$i$, $A_i$ its amplitude, and the summation extends over a band optimized per instrument (typically \SIrange{600}{1400}{\hertz} or \SIrange{1000}{2000}{\hertz}).

\begin{table}[!t]
\centering
\caption{Stability Comparison: $\Fp$ vs.\ $\Fn{2}$ (Ordinario)}
\label{tab:fp_stability}
\begin{tabular}{@{}lS[table-format=4.0]S[table-format=3.0]S[table-format=4.0]S[table-format=4.0]c@{}}
\toprule
Instrument & {$\Fp$ (Hz)} & {$\sigma_{\Fp}$} & {$\Fn{2}$ (Hz)} & {$\sigma_{\Fn{2}}$} & {Gain} \\
\midrule
Violin     & 893  & 92  & 1518 & 651  & $7.1\times$ \\
Trumpet    & 1046 & 98  & 1324 & 1018 & $10.4\times$ \\
Cl.\ B$\flat$ & 1412 & 169 & 1130 & 795  & $4.7\times$ \\
Horn       & 738  & 112 & 1106 & 734  & $6.5\times$ \\
\bottomrule
\end{tabular}
\end{table}

Table~\ref{tab:fp_stability} demonstrates that $\Fp$ is systematically more stable than either $\Fn{1}$ or $\Fn{2}$, with stability gains ranging from $4.7\times$ to $10.4\times$. The $\Fp$ corresponds to the \emph{Hauptformant} described by Meyer~\cite{meyer2009} for strings and represents the timbral coloration perceived by the listener regardless of the note being played. A total of 22~out of 30~instruments benefit from a validated $\Fp$ measurement.

Furthermore, $\sigma(\Fp)$ serves as a proxy for spectral bandwidth: a small $\sigma(\Fp)$ indicates a narrow, well-defined resonance (high selectivity $Q$), while a large $\sigma(\Fp)$ indicates a broad, diffuse formant. Thus, the violin ($\sigma = \SI{92}{\hertz}$) has a spectrally more focused timbre than the B$\flat$~clarinet ($\sigma = \SI{169}{\hertz}$), consistent with perceptual experience.

\subsection{Vowel Correspondence Framework}

Following Meyer~\cite{meyer2009}, we map formant zones to the French/German vowel system, extended to the full corpus:

\begin{table}[!t]
\centering
\caption{Vowel--Frequency Correspondence (after Meyer 2009)}
\label{tab:vowels}
\begin{tabular}{@{}clll@{}}
\toprule
Vowel & Range (Hz) & Timbre quality & Representative instr. \\
\midrule
/u/ & 200--400   & Depth, gravity       & Cb (172), Tuba (226) \\
/o/ & 400--600   & Fullness, roundness  & Hn (388), Bsn (495) \\
/\aa/ & 600--800 & Warmth, transition   & Ob (743), Tpt (786) \\
/a/ & 800--1250  & Power, presence      & Vln $\Fp$ (893) \\
/e/ & 1250--2600 & Clarity, brilliance  & Fl $\Fp$ (1535) \\
/i/ & $>$2600    & Extreme brilliance   & Tpt+harmon (2358) \\
\bottomrule
\end{tabular}
\end{table}

% ============================================================
\section{Results by Instrument Family}
\label{sec:results}

\begin{figure*}[!t]
\centering
\includegraphics[width=\textwidth]{synthese_fig1_positions.png}
\caption{Formant positions ($\Fn{1}$--$\Fn{4}$) of all 27~instruments on a logarithmic frequency axis (\SIrange{100}{5000}{\hertz}). Marker size and opacity decrease from $\Fn{1}$ (largest) to $\Fn{4}$ (smallest), reflecting decreasing perceptual salience. Green diamonds ($\blacklozenge$) indicate $\Fp$ centroid values where computed. Colored bands correspond to Meyer's vowel zones.}
\label{fig:positions}
\end{figure*}

\subsection{Woodwinds}

The woodwind family spans the widest formant range of any orchestral family (\SIrange{226}{2638}{\hertz}, vowels /u/$\to$/i/), reflecting the diversity of bore profiles and excitation mechanisms (Fig.~\ref{fig:positions}).

\subsubsection{Double Reeds}
The oboe ($\Fn{1} = \SI{743}{\hertz}$, zone /\aa/), English horn ($\Fn{1} = \SI{452}{\hertz}$, zone /o/), bassoon ($\Fn{1} = \SI{495}{\hertz}$, zone /o/), and contrabassoon ($\Fn{1} = \SI{226}{\hertz}$, zone /u/) share a common energy profile around \SIrange{900}{1200}{\hertz} via their $\Fp$ values, explaining their natural tendency to blend in the orchestra. The bassoon achieves \emph{unanimous agreement} across all four academic sources: Backus (\SIrange{440}{500}{\hertz}), Giesler (\SI{500}{\hertz}), Meyer ($\sim$\SI{500}{\hertz}), and SOL2020 (\SI{495}{\hertz}).

\subsubsection{Clarinets}
The B$\flat$~clarinet (cylindrical closed bore) presents no fixed formant in the strict sense: the spectral peak follows the fundamental (ratio~1.0) or the third harmonic (ratio~3.0), consistent with the dominance of odd partials. $\Fn{1} = \SI{463}{\hertz}$ (first cylindrical tube mode), while $\Fp = \SI{1412}{\hertz}$ is $4.7\times$ more stable than $\Fn{2}$. Divergence among sources reflects register-dependent measurements: Backus and Meyer (\SI{1500}{\hertz}) capture the \emph{clairon} register; Giesler (\SIrange{3000}{4000}{\hertz}) the extreme upper register.

\subsubsection{Flutes}
Flutes have no fixed formant (confirmed by Meyer and Giesler): the spectrum decays quasi-linearly above the fundamental. $\Fn{1} = \SI{743}{\hertz}$ (zone /\aa/) with $\sigma = \SI{324}{\hertz}$, while $\Fp = \SI{1535}{\hertz}$ provides more stable characterization. The flute--oboe pair achieves $\Delta = \SI{0}{\hertz}$---perfect formant unison between two different families.

\subsection{Brass}

\subsubsection{Open Instruments}
The horn ($\Fn{1} = \SI{388}{\hertz}$, zone /o/) achieves unanimous agreement across all academic sources and converges with the viola at $\Delta = \SI{11}{\hertz}$. The trumpet ($\Fn{1} = \SI{786}{\hertz}$, zone /\aa/) is the most register-variable brass instrument ($\sigma_{\Fn{1}} = \SI{886}{\hertz}$), but its $\Fp = \SI{1046}{\hertz}$ ($\sigma = \SI{98}{\hertz}$) is $10.4\times$ more stable. The trombone ($\Fn{1} = \SI{237}{\hertz}$, zone /u/) and tuba ($\Fn{1} = \SI{226}{\hertz}$, zone /u/) form a /u/ cluster at $\Delta = \SI{11}{\hertz}$, joined by the contrabassoon ($\Fn{1} = \SI{226}{\hertz}$, $\Delta = \SI{0}{\hertz}$).

\subsubsection{Muted Brass}
Mutes produce highly non-uniform effects depending on mute type and instrument:
\begin{itemize}
\item \textbf{Strings and horn:} systematic $\Fn{1}$ lowering (6--40\%), consistent with naive expectation.
\item \textbf{Cup mute:} $\Fn{1}$ \emph{rises} (trombone $+54\%$, trumpet $+84\%$)---the mute filters lows and pushes energy upward.
\item \textbf{Harmon mute:} extreme and instrument-dependent: trumpet $\Fn{1}$ rises from \SI{786}{\hertz} to \SI{2358}{\hertz} ($+200\%$, zone /i/); trombone $\Fn{1}$ drops to \SI{162}{\hertz} ($-32\%$).
\end{itemize}
The harmon trumpet represents the most extreme timbral transposition in the corpus: a two-vowel upward shift from /a/ to /i/.

\subsection{Strings}

\subsubsection{Solo Instruments}
The violin ($\Fn{1} = \SI{506}{\hertz}$, zone /o/) converges with the bassoon at $\Delta = \SI{11}{\hertz}$---the tightest cross-family convergence in the corpus after the flute--oboe unison. Its $\Fp = \SI{893}{\hertz}$ ($\sigma = \SI{92}{\hertz}$) corresponds to Meyer's \emph{Hauptformant}. The viola ($\Fn{1} = \SI{377}{\hertz}$, zone /\aa/--/o/) converges with the horn at $\Delta = \SI{11}{\hertz}$. The cello ($\Fn{1} = \SI{205}{\hertz}$, zone /u/) converges with the trombone ($\Delta = \SI{32}{\hertz}$) and the tuba ($\Delta = \SI{21}{\hertz}$). The double bass ($\Fn{1} = \SI{172}{\hertz}$) anchors the orchestral foundation in deep /u/.

\subsubsection{Ensemble Section Effect}
Comparison of solo vs.\ ensemble data reveals that the section effect does \emph{not} shift $\Fn{1}$ upward as might be expected, but rather \emph{flattens} $\Fn{2}$ and $\Fn{3}$:

\begin{table}[!t]
\centering
\caption{Solo vs.\ Ensemble Section Effect}
\label{tab:section}
\begin{tabular}{@{}lS[table-format=3.0]S[table-format=3.0]cS[table-format=4.0]S[table-format=4.0]c@{}}
\toprule
 & {Solo $\Fn{1}$} & {Ens.\ $\Fn{1}$} & {$\Delta\Fn{1}$} & {Solo $\Fn{2}$} & {Ens.\ $\Fn{2}$} & {$\Delta\Fn{2}$} \\
\midrule
Violin & 506  & 495  & $-2\%$  & 1518 & 1163 & $-23\%$ \\
Viola  & 377  & 366  & $-3\%$  & 829  & 764  & $-8\%$  \\
Cello  & 205  & 205  & $0\%$   & 506  & 474  & $-6\%$  \\
D.~bass & 172 & 172  & $0\%$   & 474  & 441  & $-7\%$  \\
\bottomrule
\end{tabular}
\end{table}

The perceptual consequence is a \emph{blended}, less individually characterized timbre---the characteristic texture difference between a string quartet and a full orchestral section.

\subsection{Saxophone}

Only the alto saxophone is represented in the SOL2020 corpus. Its $\Fn{1} = \SI{398}{\hertz}$ (zone /o/) places it squarely within the convergence cluster, at $\Delta = \SI{10}{\hertz}$ from the horn---explaining the well-known timbral affinity between these instruments. Giesler~\cite{giesler1985} explicitly groups the tenor saxophone with bassoon, horn, and trombone in the /o/ zone.

% ============================================================
\section{Timbral Convergence Analysis}
\label{sec:convergence}

\subsection{The /o/--/\aa/ Convergence Cluster}

The central finding of this study is that the vowel zone /o/--/\aa/ (\SIrange{377}{506}{\hertz}) acts as a multi-family convergence area (Figs.~\ref{fig:vowelspace}--\ref{fig:cluster}). Within a span of only \SI{129}{\hertz}, 11~instruments from four families share nearly identical $\Fn{1}$ values:

\begin{table}[!t]
\centering
\caption{The /o/--/\aa/ Convergence Cluster}
\label{tab:cluster}
\begin{tabular}{@{}llS[table-format=3.0]cS[table-format=3.0]@{}}
\toprule
Instrument & Family & {$\Fn{1}$ (Hz)} & Vowel & {$\Delta$ vs.\ Vln} \\
\midrule
Viola             & Strings    & 377 & /o/ & 129 \\
Horn              & Brass      & 388 & /o/ & 118 \\
Alto sax          & Saxophones & 398 & /o/ & 108 \\
English horn      & Woodwinds  & 452 & /o/ & 54 \\
Cl.\ B$\flat$     & Woodwinds  & 463 & /o/ & 43 \\
Bassoon           & Woodwinds  & 495 & /o/ & 11 \\
Vln.\ ensemble    & Strings    & 495 & /o/ & 11 \\
Violin            & Strings    & 506 & /o/ & {---} \\
\bottomrule
\end{tabular}
\end{table}

\begin{figure}[!t]
\centering
\includegraphics[width=\columnwidth]{synthese_fig2_espace_vocalique.png}
\caption{$\Fn{1}$/$\Fn{2}$ vowel space of 27~instruments (phonetic convention: $\Fn{1}$ horizontal, $\Fn{2}$ vertical). Marker shapes distinguish families: $\blacklozenge$~Woodwinds, $\blacksquare$~Brass, $\bullet$~Solo strings, $\blacktriangle$~Ensemble strings, $+$~Saxophones. The dashed red ellipse delineates the /o/--/\aa/ convergence zone ($\Fn{1} = \SIrange{377}{506}{\hertz}$).}
\label{fig:vowelspace}
\end{figure}

\begin{figure}[!t]
\centering
\includegraphics[width=\columnwidth]{synthese_fig3_cluster.png}
\caption{Schematic spectral envelopes (Gaussian profiles centered on $\Fn{1}$) of the 11~instruments converging in the /o/--/\aa/ zone (\SIrange{377}{506}{\hertz}, red band). Overlapping resonances in a narrow band produce the perception of a single, homogeneous timbre---the acoustic mechanism of orchestral fusion.}
\label{fig:cluster}
\end{figure}

\subsection{Verified Doublings}

Table~\ref{tab:doublings} presents the 20~most significant doublings, ranked by $\Delta$ (ascending). These doublings span four quality tiers:
\begin{itemize}
\item $\Delta = \SI{0}{\hertz}$: formant unison (3~pairs);
\item $\Delta \leq \SI{30}{\hertz}$: quasi-perfect fusion (7~pairs);
\item $\Delta$ \SIrange{30}{80}{\hertz}: excellent convergence (2~pairs);
\item $\Delta > \SI{200}{\hertz}$: spectral complementarity (8~pairs).
\end{itemize}

\begin{table*}[!t]
\centering
\caption{Top 12 Verified Doublings (by Ascending $\Delta$)}
\label{tab:doublings}
\begin{tabular}{@{}clS[table-format=4.0]lS[table-format=4.0]S[table-format=3.0]ll@{}}
\toprule
{\#} & Instrument A & {$\Fn{1}$ A} & Instrument B & {$\Fn{1}$ B} & {$\Delta$ (Hz)} & Quality & Zone \\
\midrule
1  & Flute           & 743  & Oboe             & 743  & 0  & Unison $\star$    & /\aa/ \\
2  & Tuba CB         & 226  & Contrabassoon    & 226  & 0  & Unison $\star$    & /u/ \\
3  & Tuba CB         & 226  & Tuba bass        & 226  & 0  & Unison $\star$    & /u/ \\
4  & Horn            & 388  & Viola            & 377  & 11 & Quasi-perfect     & /o/ \\
5  & English hn.     & 452  & Cl.\ B$\flat$    & 463  & 11 & Quasi-perfect     & /o/ \\
6  & Violin          & 506  & Bassoon          & 495  & 11 & Quasi-perfect     & /o/ \\
7  & Trombone        & 237  & Tuba bass        & 226  & 11 & Quasi-perfect     & /u/ \\
8  & Trombone        & 237  & Contrabassoon    & 226  & 11 & Quasi-perfect     & /u/ \\
9  & Trombone        & 237  & Trb.\ bass       & 258  & 21 & Quasi-perfect     & /u/ \\
10 & Contrabassoon   & 226  & Double bass      & 172  & 54 & Excellent         & /u/ \\
11 & English hn.     & 452  & Horn             & 388  & 64 & Excellent         & /o/ \\
12 & Oboe            & 743  & Bassoon          & 495  & 248& Complementary     & /\aa/+/o/ \\
\bottomrule
\end{tabular}
\end{table*}

\begin{figure}[!t]
\centering
\includegraphics[width=\columnwidth]{synthese_f1_base.png}
\caption{$\Fn{1}$ and $\Fp$ positions for the 20~base instruments (ordinario), sorted by ascending $\Fn{1}$. Vowel zones are indicated by colored background bands. The /o/--/\aa/ convergence zone (\SIrange{377}{506}{\hertz}) concentrates the highest density of instruments.}
\label{fig:f1base}
\end{figure}

\subsection{Convergence Matrix}

The full convergence matrix (Fig.~\ref{fig:matrix}) plots pairwise $\Delta\Fn{1}$ values for all 20~base instruments. Green cells ($\Delta \leq \SI{80}{\hertz}$) identify fusional doublings; yellow cells (\SIrange{80}{200}{\hertz}) indicate masking risk zones; red cells ($\Delta > \SI{200}{\hertz}$) indicate spectral complementarity.

\begin{figure}[!t]
\centering
\includegraphics[width=\columnwidth]{synthese_matrix_base.png}
\caption{Pairwise $\Fn{1}$ convergence matrix for 20~base instruments. Green: $\Delta \leq \SI{80}{\hertz}$ (fusion); yellow: \SIrange{80}{200}{\hertz} (masking risk); red: $\Delta > \SI{200}{\hertz}$ (complementarity).}
\label{fig:matrix}
\end{figure}

% ============================================================
\section{Six Principles of Acoustic Orchestration}
\label{sec:principles}

The data support six quantitative principles:

\subsection{Principle 1: Formant Convergence (Fusion)}
When two instruments share a similar $\Fn{1}$ ($\Delta \leq \SI{80}{\hertz}$), the ear perceives a single homogeneous timbre rather than two distinct voices. This fusion operates regardless of the melodic interval---bassoon and violin fuse at \emph{any} harmonic interval because formant convergence is a \emph{timbral} property, not a pitch property.

\subsection{Principle 2: Spectral Complementarity (Enrichment)}
When $\Delta\Fn{1} > \SI{200}{\hertz}$, the two spectra complement each other without mutual interference. Optimal enrichment occurs across at least two vowel zones ($\Delta \geq \SIrange{200}{300}{\hertz}$). Example: trumpet (\SI{786}{\hertz}, /\aa/) + trombone (\SI{237}{\hertz}, /u/), $\Delta = \SI{549}{\hertz}$.

\subsection{Principle 3: Ensemble Section Effect}
The transition from soloist to ensemble does not shift $\Fn{1}$ but flattens $\Fn{2}$/$\Fn{3}$ by 6--23\% (Table~\ref{tab:section}). Individual spectral irregularities average out in the collective, producing a \emph{blended} texture with reduced spectral relief.

\subsection{Principle 4: Mute as Timbral Transposition}
For strings and horn, mutes systematically lower $\Fn{1}$. For brass cup and harmon mutes, the effect is \emph{reversed}---$\Fn{1}$ rises, sometimes dramatically (trumpet harmon: $+200\%$). The mute is thus an unpredictable timbral transposition tool, not a simple spectral attenuator.

\subsection{Principle 5: Cross-Family Formant Groups}
Traditional orchestral families (woodwinds, brass, strings) do not align with formant-based groupings. The analysis by vowel zones reveals four \emph{cross-family acoustic groupings} that traverse organological boundaries (Fig.~\ref{fig:f1base}):
\begin{itemize}
\item \textbf{/u/ family} (\SIrange{170}{330}{\hertz}): double bass, cello, trombone, tuba, contrabassoon, bass clarinet.
\item \textbf{/o/--/\aa/ family} (\SIrange{377}{510}{\hertz}): viola, horn, alto sax, English horn, B$\flat$~clarinet, bassoon, violin.
\item \textbf{/a/ family} (\SIrange{800}{1250}{\hertz} via $\Fp$): trumpet, horn, English horn, bass clarinet, bass trombone.
\item \textbf{/e/ family} (\SIrange{1250}{2600}{\hertz} via $\Fp$): flute, oboe, E$\flat$~clarinet.
\end{itemize}
These cross-family groups are more predictive of timbral fusion than organological categories.

\subsection{Principle 6: Spectral Masking Avoidance}
The highest masking risk occurs at $\Delta\Fn{1}$ \SIrange{80}{200}{\hertz}---partial overlap without fusion. Identified risk pairs include B$\flat$~clarinet--viola ($\Delta = \SI{86}{\hertz}$), bassoon--horn ($\Delta = \SI{107}{\hertz}$), and horn--violin ($\Delta = \SI{118}{\hertz}$). Three avoidance strategies: push toward convergence ($\Delta \leq \SI{30}{\hertz}$), push toward divergence ($\Delta \geq \SI{300}{\hertz}$), or differentiate dynamics/register.

% ============================================================
\section{Validation}
\label{sec:validation}

\subsection{Pipeline Integrity}
The extraction pipeline v22 reproduces all 16~SOL2020 instruments at $\Delta = \SI{0}{\hertz}$ versus raw spectral envelope data. This end-to-end audit eliminates pipeline artifacts as a source of error.

\subsection{Multi-Source Agreement}
Cross-validation against Backus, Giesler, Meyer, and McCarty yields a global agreement rate of \textbf{93\%} (27/29 comparable instruments). Perfect unanimity is achieved for the bassoon (four sources agree within \SI{60}{\hertz}). Apparent disagreements for the trombone and cello are resolved by recognizing that academic sources often measured $\Fp$ or upper-register peaks rather than the strict $\Fn{1}$.

\subsection{Addressing Audit Observations}
Two independent audits raised methodological concerns that we address here:
\begin{enumerate}
\item \textbf{$\Fn{1}$/$\Fp$ confusion in reference tables:} In earlier versions of this work, some tables mixed $\Fp$ centroid values in $\Fn{1}$ columns. The current version strictly separates $\Fn{1}$--$\Fn{6}$ (peak-based) from $\Fp$ (centroid-based) with explicit column headers. $\Fp$ values are always marked with the $\Fp$ symbol.
\item \textbf{Violin ensemble $\Fn{1}$ statement:} An earlier synthesis section incorrectly stated that the violin ensemble $\Fn{1}$ was \SI{1556}{\hertz}. The correct value is \SI{495}{\hertz} ($-2\%$ vs.\ solo); the section effect manifests primarily in $\Fn{2}$ ($-23\%$) and $\Fn{3}$ ($-16\%$), not in $\Fn{1}$.
\item \textbf{1D vs.\ 2D vowel mapping:} The vowel correspondence table maps instruments via a single frequency ($\Fn{1}$), whereas phonetic vowel identification requires at minimum the ($\Fn{1}$, $\Fn{2}$) plane. We acknowledge this limitation while noting that for orchestration purposes, the $\Fn{1}$-based convergence metric is the primary predictor of timbral fusion; a full ($\Fn{1}$, $\Fn{2}$) analysis (Fig.~\ref{fig:vowelspace}) is provided in the synthesis.
\item \textbf{Bandwidth characterization:} Rather than measuring formant bandwidth ($Q$ factor) directly---which requires parametric curve fitting sensitive to model assumptions---we use $\sigma(\Fp)$ as a statistically robust indicator of spectral selectivity (Section~\ref{sec:fp}). A low $\sigma(\Fp)$ corresponds to a narrow, well-defined resonance (high $Q$); a high $\sigma(\Fp)$ to a broad, diffuse one. This approach captures the same perceptual distinction while being computed directly from the data without additional modeling.
\end{enumerate}

% ============================================================
\section{Discussion}
\label{sec:discussion}

\subsection{Limitations}
Several limitations should be noted:
\begin{itemize}
\item The saxophone family is represented by the alto only; soprano, tenor, and baritone are absent from the corpus.
\item The study analyzes sustained tones only; formant transitions during attacks and dynamic changes---critical for timbral perception in musical context---are not captured.
\item The vowel correspondence framework maps zones via $\Fn{1}$ alone. A full ($\Fn{1}$, $\Fn{2}$) mapping on a Bark or ERB scale would strengthen the phonetic analogy.
\item Robust statistics (median, IQR, 95\% confidence intervals) are reported for individual instruments but not systematically for all inter-instrument comparisons.
\end{itemize}

\subsection{Implications for Compositional Practice}
The formant convergence framework provides a quantitative basis for orchestration decisions that have historically relied on intuition and tradition. A composer seeking to reinforce the /o/ color can freely double among English horn, horn, bassoon, and violin---the acoustic family transcends instrumental families. Conversely, combining two formant families produces the spectral enrichment of Principle~2.

The $\Fp$ metric opens new possibilities for analyzing spectral orchestration in contemporary repertoire (Grisey, Murail, Saariaho, Haas, Radulescu), where timbral blending is an explicit compositional parameter. Future work will apply this methodology to extract doublings from scores and validate them against spectral analysis of recordings.

% ============================================================
\section{Conclusion}
\label{sec:conclusion}

This study provides the first comprehensive, corpus-validated quantitative reference for orchestral instrument formants, covering 30~instruments and \num{5914}~sustained samples. The introduction of the energy-weighted spectral centroid $\Fp$ resolves the long-standing instability of peak-based formant extraction for wide-tessitura instruments, offering stability gains of up to $10.4\times$.

The discovery of the /o/--/\aa/ convergence cluster---where 11~instruments from four families converge within \SI{129}{\hertz}---provides a measurable acoustic explanation for classical orchestral doublings. The six derived orchestration principles offer a quantitative framework that bridges acoustic science and compositional practice, validated at 93\% against four independent academic sources.

Future extensions include expanding the corpus to the full saxophone family and percussion, mapping the ($\Fn{1}$, $\Fn{2}$) vowel space on perceptual scales, and applying the convergence framework to the analysis of spectral orchestration in contemporary music.

% ============================================================
\section*{Acknowledgment}
The author thanks IRCAM for the SOL2020 database and the Ensemble Intercontemporain for recording access.

\balance
\begin{thebibliography}{10}

\bibitem{fant1960}
G.~Fant, \emph{Acoustic Theory of Speech Production}.\hskip 1em plus 0.5em minus 0.4em The Hague: Mouton, 1960.

\bibitem{backus1969}
J.~Backus, \emph{The Acoustical Foundations of Music}.\hskip 1em plus 0.5em minus 0.4em New York: W.\,W.~Norton, 1969.

\bibitem{giesler1985}
W.~Gieseler, L.~Lombardi, and R.\,D.~Weyer, \emph{Instrumentation in der Musik des 20.\ Jahrhunderts: Akustik, Instrumente, Zusammenwirken}.\hskip 1em plus 0.5em minus 0.4em Celle: Moeck, 1985.

\bibitem{meyer2009}
J.~Meyer, \emph{Acoustics and the Performance of Music}, 5th~ed.\hskip 1em plus 0.5em minus 0.4em New York: Springer, 2009.

\bibitem{mccarty2003}
J.~McCarty, ``Formant analysis of orchestral instruments,'' Center for Computer Research in Music and Acoustics (CCRMA), Stanford Univ., Stanford, CA, 2003. [Online]. Available: \url{https://ccrma.stanford.edu/~jmccarty/formant.htm}

\bibitem{benade1990}
A.\,H.~Benade, \emph{Fundamentals of Musical Acoustics}, 2nd~ed.\hskip 1em plus 0.5em minus 0.4em New York: Dover, 1990.

\bibitem{sol2020}
IRCAM, ``SOL---Studio On Line: Instrumental sounds datasets,'' 2020. [Online]. Available: \url{https://forum.ircam.fr/collections/detail/sol-instrumental-sounds-datasets/}

\bibitem{orchideasol}
C.~Music, ``OrchideaSOL,'' Zenodo, 2020. [Online]. Available: \url{https://zenodo.org/records/3686252}

\bibitem{grey1977}
J.\,M.~Grey, ``Multidimensional perceptual scaling of musical timbres,'' \emph{J.\ Acoust.\ Soc.\ Am.}, vol.~61, no.~5, pp.~1270--1277, 1977.

\bibitem{peeters2011}
G.~Peeters, B.\,L.~Giordano, P.~Susini, N.~Misdariis, and S.~McAdams, ``The timbre toolbox: Extracting audio descriptors from musical signals,'' \emph{J.\ Acoust.\ Soc.\ Am.}, vol.~130, no.~5, pp.~2902--2916, 2011.

\end{thebibliography}

\end{document}
